% In this file you should put the actual content of the blueprint.
% It will be used both by the web and the print version.
% It should *not* include the \begin{document}
%
% If you want to split the blueprint content into several files then
% the current file can be a simple sequence of \input. Otherwise It
% can start with a \section or \chapter for instance.
\begin{definition}\label{def:formula}
  \lean{Formula}
  \leanok
  The set of propositional formulas is defined inductively:
  $\bot$ (falsum), propositional variables $p_n$ ($n \in \mathbb{N}$), and closed
  under negation, conjunction, disjunction, and implication.
\end{definition}
\begin{definition}\label{def:formula_tostr}
  \lean{Formula.toStr}
  \leanok
  \uses{def:formula}
  Pretty-printing a formula as a string.
\end{definition}
\begin{definition}\label{def:formula_tonat}
  \lean{Formula.toNat}
  \leanok
  \uses{def:formula}
  Encodes a formula as a natural number via Cantor pairing.
\end{definition}
\begin{definition}\label{def:formula_ofnat}
  \lean{Formula.ofNat}
  \leanok
  \uses{def:formula}
  Decodes a natural number back into an (optional) formula.
\end{definition}

\begin{theorem}\label{thm:formula_ofnat_tonat}
  \lean{Formula.ofNat_toNat}
  \leanok
  \uses{def:formula_tonat, def:formula_ofnat}
  Decoding after encoding recovers the original formula:
  $\mathrm{ofNat}(\mathrm{toNat}(\varphi)) = \varphi$.
\end{theorem}

\begin{definition}\label{def:formula_encodable}
  \lean{instEncodableFormula}
  \leanok
  \uses{def:formula_tonat, def:formula_ofnat, thm:formula_ofnat_tonat}
  The \texttt{Encodable} instance for \texttt{Formula}, giving an effective
  enumeration of all formulas.
\end{definition}

\begin{definition}\label{def:formula_syntactic_identity}
  \lean{Formula.syntatic_identity}
  \leanok
  \uses{def:formula_tostr}
  Two formulas are syntactically identical iff their string representations are equal.
\end{definition}

\begin{definition}\label{def:formula_complexity}
  \lean{Formula.complexity}
  \leanok
  \uses{def:formula}
  The structural complexity of a formula, used for induction.
\end{definition}

\begin{definition}\label{def:formula_lp}
  \lean{Formula.LP}
  \leanok
  \uses{def:formula_tostr}
  Counts left parentheses in a formula's string representation.
\end{definition}

\begin{definition}\label{def:formula_rp}
  \lean{Formula.RP}
  \leanok
  \uses{def:formula_tostr}
  Counts right parentheses in a formula's string representation.
\end{definition}

\begin{theorem}\label{thm:formula_lp_count_monotonicity}
  \lean{Formula.LP.count_monotonicity}
  \leanok
  \uses{def:formula_lp}
  Left-parenthesis count distributes over list cons:
  $\mathrm{count}(h::t) = \mathrm{count}([h]) + \mathrm{count}(t)$.
\end{theorem}

\begin{theorem}\label{thm:formula_rp_count_monotonicity}
  \lean{Formula.RP.count_monotonicity}
  \leanok
  \uses{def:formula_rp}
  Right-parenthesis count distributes over list cons.
\end{theorem}

\begin{theorem}\label{thm:formula_lp_count_append_monotonicity}
  \lean{Formula.LP.count_append_monotonicity}
  \leanok
  \uses{thm:formula_lp_count_monotonicity}
  Left-parenthesis count distributes over list append:
  $\mathrm{count}(h \mathbin{+\!\!+} t) = \mathrm{count}(h) + \mathrm{count}(t)$.
\end{theorem}

\begin{theorem}\label{thm:formula_rp_count_append_monotonicity}
  \lean{Formula.RP.count_append_monotonicity}
  \leanok
  \uses{thm:formula_rp_count_monotonicity}
  Right-parenthesis count distributes over list append.
\end{theorem}

\begin{definition}\label{def:formula_uniformsub}
  \lean{Formula.uniformSub}
  \leanok
  \uses{def:formula}
  Uniform substitution: replaces every occurrence of variable $n$ by variable $n'$
  throughout a formula.
\end{definition}

\begin{axiom}\label{ax:dne_elim}
  \lean{dne_elim}
  \leanok
  \uses{def:formula}
  Double negation identity: $\varphi = \neg\neg\varphi$, assumed as an axiom
  rather than derived syntactically.
\end{axiom}
